% =====================================================================
\chapter{Sucesiones y límites}
\label{tema:2}
% =====================================================================

El concepto central del análisis es el de límite. Este tema lo introduce primero para sucesiones, donde la definición es más sencilla, y después para funciones, con la definición $\varepsilon$--$\delta$. Se demuestran las reglas de cálculo de límites, el teorema del sándwich y el teorema de la convergencia monótona; se obtienen los límites notables que se usarán en el resto del módulo y se estudian las asíntotas. Se siguen los capítulos~5 y~22 de Spivak~\cite{spivak}, el capítulo~2 y la §5.1 de los apuntes de \emph{Analysis~I} de Cambridge~\cite{cambA1}, y las §2.2--2.5 y~4.6 de OpenStax~\cite{osc1}.

\section{Sucesiones convergentes}

\begin{definicion}[Sucesión y límite]\label{def:limite-sucesion}
Una \emph{sucesión} de números reales es una función $\N\to\R$, $n\mapsto a_n$; se escribe $(a_n)_{n\ge1}$ o $(a_n)$. Se dice que $(a_n)$ \emph{converge} a $L\in\R$, y se escribe $a_n\to L$ o $\lim_{n\to\infty}a_n=L$, si
\[
\forall\varepsilon>0\ \ \exists N\in\N\ \ \forall n\ge N:\quad \abs{a_n-L}<\varepsilon.
\]
Una sucesión que no converge a ningún número real es \emph{divergente}.
\end{definicion}

La definición se lee como un juego: un adversario fija una tolerancia $\varepsilon>0$, tan pequeña como quiera, y hay que responder con un índice $N$ a partir del cual todos los términos están a distancia menor que $\varepsilon$ de $L$. El índice $N$ puede depender de $\varepsilon$, y casi siempre depende.

\begin{ejemplo}[$1/n\to0$]\label{ej:1n}
Sea $\varepsilon>0$. Por la propiedad arquimediana de $\R$ (módulo~0.3) existe $N\in\N$ con $N>1/\varepsilon$. Para $n\ge N$, $\abs{1/n-0}=1/n\le1/N<\varepsilon$. Luego $1/n\to0$.
\end{ejemplo}

\begin{ejemplo}[Una demostración $\varepsilon$--$N$ con acotación previa]\label{ej:eN}
Se demuestra que $a_n=\frac{3n+1}{n+2}\to3$. \emph{Trabajo previo:}
\[
\abs{a_n-3}=\Bigl|\frac{3n+1-3n-6}{n+2}\Bigr|=\frac{5}{n+2}<\frac5n.
\]
Basta, pues, $5/n<\varepsilon$, es decir, $n>5/\varepsilon$. \emph{Demostración:} dado $\varepsilon>0$, sea $N\in\N$ con $N>5/\varepsilon$. Para $n\ge N$, $\abs{a_n-3}<5/n\le5/N<\varepsilon$. \hfill$\blacksquare$
\end{ejemplo}

En las demostraciones $\varepsilon$--$N$ conviene separar el trabajo previo (acotar $\abs{a_n-L}$ por una expresión sencilla que tienda a~0) de la demostración en sí, que se redacta en el orden de los cuantificadores: «sea $\varepsilon>0$; sea $N=\dots$; para $n\ge N$, \dots».

\begin{proposicion}[Unicidad y acotación]\label{prop:unicidad}
(i) El límite de una sucesión convergente es único. (ii) Toda sucesión convergente está acotada.
\end{proposicion}

\begin{proof}
(i) Supóngase $a_n\to L$ y $a_n\to M$ con $L\ne M$, y sea $\varepsilon=\abs{L-M}/2>0$. Existen $N_1$ y $N_2$ tales que $\abs{a_n-L}<\varepsilon$ para $n\ge N_1$ y $\abs{a_n-M}<\varepsilon$ para $n\ge N_2$. Para $n=\max(N_1,N_2)$, la desigualdad triangular da $\abs{L-M}\le\abs{L-a_n}+\abs{a_n-M}<2\varepsilon=\abs{L-M}$, contradicción.
(ii) Sea $a_n\to L$. Con $\varepsilon=1$ existe $N$ tal que $\abs{a_n}\le\abs{a_n-L}+\abs{L}<1+\abs{L}$ para $n\ge N$. Entonces $\abs{a_n}\le K=\max\bigl(\abs{a_1},\dots,\abs{a_{N-1}},1+\abs{L}\bigr)$ para todo $n$.
\end{proof}

\begin{teorema}[Álgebra de límites de sucesiones]\label{teo:algebra-sucesiones}
Sean $a_n\to L$ y $b_n\to M$. Entonces
(i) $a_n+b_n\to L+M$; (ii) $ca_n\to cL$ para todo $c\in\R$; (iii) $a_nb_n\to LM$; (iv) si $M\ne0$ y $b_n\ne0$ para todo $n$, $a_n/b_n\to L/M$.
\end{teorema}

\begin{proof}
(i) Dado $\varepsilon>0$, existen $N_1$, $N_2$ con $\abs{a_n-L}<\varepsilon/2$ para $n\ge N_1$ y $\abs{b_n-M}<\varepsilon/2$ para $n\ge N_2$. Para $n\ge\max(N_1,N_2)$, $\abs{(a_n+b_n)-(L+M)}\le\abs{a_n-L}+\abs{b_n-M}<\varepsilon$.
(iii) Se escribe $a_nb_n-LM=a_n(b_n-M)+M(a_n-L)$. Por la \cref{prop:unicidad}, $\abs{a_n}\le K$ para algún $K>0$. Dado $\varepsilon>0$, existen $N_1,N_2$ con $\abs{b_n-M}<\varepsilon/(2K)$ para $n\ge N_1$ y $\abs{a_n-L}<\varepsilon/(2\abs{M}+1)$ para $n\ge N_2$. Para $n\ge\max(N_1,N_2)$, $\abs{a_nb_n-LM}\le K\abs{b_n-M}+\abs{M}\abs{a_n-L}<\varepsilon/2+\varepsilon/2$.
(ii) Es (iii) con la sucesión constante $b_n=c$. (iv) Basta ver que $1/b_n\to1/M$ y aplicar (iii); es el \cref{eje:2-cociente}.
\end{proof}

\begin{teorema}[Teorema del sándwich]\label{teo:sandwich}
Si $a_n\le b_n\le c_n$ para todo $n\ge n_0$ y $a_n\to L$, $c_n\to L$, entonces $b_n\to L$.
\end{teorema}

\begin{proof}
Dado $\varepsilon>0$, para $n$ suficientemente grande se tiene $L-\varepsilon<a_n$ y $c_n<L+\varepsilon$, luego $L-\varepsilon<a_n\le b_n\le c_n<L+\varepsilon$, es decir, $\abs{b_n-L}<\varepsilon$.
\end{proof}

\begin{definicion}[Divergencia a infinito]\label{def:div-infinito}
$a_n\to+\infty$ si para todo $K\in\R$ existe $N$ tal que $a_n>K$ para todo $n\ge N$. Análogamente, $a_n\to-\infty$.
\end{definicion}

\begin{proposicion}[Sucesión geométrica]\label{prop:geometrica}
Sea $r\in\R$. Si $\abs{r}<1$, $r^{n}\to0$; si $r>1$, $r^{n}\to+\infty$; si $r=1$, $r^{n}\to1$; si $r\le-1$, $(r^{n})$ diverge.
\end{proposicion}

\begin{proof}
Si $r>1$, sea $h=r-1>0$. Por la desigualdad de Bernoulli (módulo~0.3), $r^{n}=(1+h)^{n}\ge1+nh$, que supera cualquier $K$ para $n$ grande. Si $0<\abs{r}<1$, entonces $1/\abs{r}>1$ y, dado $\varepsilon>0$, existe $N$ con $(1/\abs{r})^{n}>1/\varepsilon$ para $n\ge N$, es decir, $\abs{r^{n}-0}<\varepsilon$. El caso $r=0$ es trivial, y $r=1$ también. Si $r\le-1$, $\abs{r^{n}}\ge1$ y los términos alternan de signo, luego dos términos consecutivos distan al menos~2 y la sucesión no puede converger (\cref{eje:2-alternada}).
\end{proof}

\section{Sucesiones monótonas y el número \texorpdfstring{$e$}{e}}

\begin{teorema}[Convergencia monótona]\label{teo:monotona}
Toda sucesión creciente y acotada superiormente converge a $\sup\{a_n:n\in\N\}$. Toda sucesión decreciente y acotada inferiormente converge a su ínfimo.
\end{teorema}

\begin{proof}
Sea $(a_n)$ creciente y acotada superiormente. Por el axioma del supremo (módulo~0.3), existe $L=\sup_n a_n$. Dado $\varepsilon>0$, $L-\varepsilon$ no es cota superior, luego existe $N$ con $a_N>L-\varepsilon$. Para $n\ge N$, el crecimiento da $L-\varepsilon<a_N\le a_n\le L$, luego $\abs{a_n-L}<\varepsilon$. El caso decreciente se obtiene aplicando lo anterior a $(-a_n)$.
\end{proof}

Este teorema es la primera consecuencia importante de la completitud de $\R$: en $\Q$ es falso, porque la sucesión creciente y acotada $1;\ 1{,}4;\ 1{,}41;\ 1{,}414;\dots$ de aproximaciones decimales de $\sqrt2$ no tiene límite racional.

\begin{teorema}[El número $e$]\label{teo:e}
La sucesión $e_n=\bigl(1+\tfrac1n\bigr)^{n}$ es creciente y está acotada superiormente por~3. Su límite se denota $e$ y vale $e=2{,}71828\dots$
\end{teorema}

\begin{proof}
Por el binomio de Newton,
\[
e_n=\sum_{k=0}^{n}\binom{n}{k}\frac{1}{n^{k}}=\sum_{k=0}^{n}\frac{1}{k!}\cdot\frac{n(n-1)\cdots(n-k+1)}{n^{k}}=\sum_{k=0}^{n}\frac{1}{k!}\prod_{j=0}^{k-1}\Bigl(1-\frac{j}{n}\Bigr).
\]
Al pasar de $n$ a $n+1$, cada factor $1-j/n$ aumenta a $1-j/(n+1)$ y se añade un término positivo más, luego $e_n<e_{n+1}$. Además, cada producto es $\le1$ y $k!\ge2^{k-1}$ para $k\ge1$, de modo que $e_n\le1+\sum_{k=1}^{n}2^{-(k-1)}<1+2=3$. Por el \cref{teo:monotona}, $(e_n)$ converge.
\end{proof}

\section{Límite de una función}

\begin{definicion}[Límite $\varepsilon$--$\delta$]\label{def:limite-funcion}
Sea $f$ definida en un intervalo abierto que contiene a $a$, salvo quizá en el propio $a$. Se dice que $f$ tiene límite $L$ en $a$, y se escribe $\lim_{x\to a}f(x)=L$, si
\[
\forall\varepsilon>0\ \ \exists\delta>0\ \ \forall x:\quad 0<\abs{x-a}<\delta\ \Rightarrow\ \abs{f(x)-L}<\varepsilon.
\]
\end{definicion}

La condición $0<\abs{x-a}$ excluye el punto $a$: el límite describe el comportamiento de $f$ \emph{cerca} de $a$, no \emph{en} $a$, y no depende de $f(a)$ ni de que esté definido.

\begin{ejemplo}[Un límite lineal]\label{ej:lineal}
$\lim_{x\to1}(2x+1)=3$. \emph{Trabajo previo:} $\abs{(2x+1)-3}=2\abs{x-1}$, que es $<\varepsilon$ si $\abs{x-1}<\varepsilon/2$. \emph{Demostración:} dado $\varepsilon>0$, sea $\delta=\varepsilon/2$. Si $0<\abs{x-1}<\delta$, entonces $\abs{(2x+1)-3}=2\abs{x-1}<2\delta=\varepsilon$. \hfill$\blacksquare$
\end{ejemplo}

\begin{ejemplo}[Un límite cuadrático: el truco del mínimo]\label{ej:cuadratico}
$\lim_{x\to2}x^{2}=4$. \emph{Trabajo previo:} $\abs{x^{2}-4}=\abs{x+2}\,\abs{x-2}$. El factor $\abs{x+2}$ no es constante, pero se acota si se restringe $x$: si $\abs{x-2}<1$, entonces $1<x<3$ y $\abs{x+2}<5$. \emph{Demostración:} dado $\varepsilon>0$, sea $\delta=\min(1,\varepsilon/5)$. Si $0<\abs{x-2}<\delta$, entonces $\abs{x-2}<1$, luego $\abs{x+2}<5$, y $\abs{x^{2}-4}<5\abs{x-2}<5\cdot\varepsilon/5=\varepsilon$. \hfill$\blacksquare$
\end{ejemplo}

\begin{teorema}[Caracterización por sucesiones]\label{teo:heine}
$\lim_{x\to a}f(x)=L$ si y solo si, para toda sucesión $(x_n)$ con $x_n\ne a$ y $x_n\to a$, se tiene $f(x_n)\to L$.
\end{teorema}

\begin{proof}
($\Rightarrow$) Sea $\varepsilon>0$ y $\delta$ el de la definición. Como $x_n\to a$, existe $N$ con $\abs{x_n-a}<\delta$ para $n\ge N$; como $x_n\ne a$, también $0<\abs{x_n-a}$, luego $\abs{f(x_n)-L}<\varepsilon$.
($\Leftarrow$) Por contrarrecíproco. Si el límite no es $L$, existe $\varepsilon>0$ tal que para todo $\delta>0$ hay un $x$ con $0<\abs{x-a}<\delta$ y $\abs{f(x)-L}\ge\varepsilon$. Tomando $\delta=1/n$ se obtiene $x_n$ con $0<\abs{x_n-a}<1/n$ y $\abs{f(x_n)-L}\ge\varepsilon$: $x_n\to a$, $x_n\ne a$ y $f(x_n)\not\to L$.
\end{proof}

El \cref{teo:heine} traslada a las funciones todo lo demostrado para sucesiones: el límite es único, y valen el álgebra de límites (\cref{teo:algebra-sucesiones}) y el teorema del sándwich para funciones. También sirve para probar que un límite \emph{no} existe: basta encontrar dos sucesiones $x_n\to a$ y $y_n\to a$ con $\lim f(x_n)\ne\lim f(y_n)$.

\begin{ejemplo}[Un límite que no existe]\label{ej:sin1x}
$\lim_{x\to0}\sin(1/x)$ no existe. Con $x_n=1/(2\pi n)$, $\sin(1/x_n)=0$; con $y_n=1/(2\pi n+\pi/2)$, $\sin(1/y_n)=1$. Ambas sucesiones tienden a~0 y las imágenes tienen límites distintos. En cambio, $\lim_{x\to0}x\sin(1/x)=0$ por el sándwich: $-\abs{x}\le x\sin(1/x)\le\abs{x}$.
\end{ejemplo}

\begin{definicion}[Límites laterales, infinitos y en el infinito]\label{def:limites-variantes}
Sea $f$ una función real y $a\in\R$.
\begin{enumerate}
  \item $\lim_{x\to a^{+}}f(x)=L$ si para todo $\varepsilon>0$ existe $\delta>0$ con $\abs{f(x)-L}<\varepsilon$ siempre que $a<x<a+\delta$. Análogamente, $\lim_{x\to a^{-}}$ con $a-\delta<x<a$.
  \item $\lim_{x\to a}f(x)=+\infty$ si para todo $K$ existe $\delta>0$ con $f(x)>K$ siempre que $0<\abs{x-a}<\delta$.
  \item $\lim_{x\to+\infty}f(x)=L$ si para todo $\varepsilon>0$ existe $M$ con $\abs{f(x)-L}<\varepsilon$ siempre que $x>M$.
\end{enumerate}
\end{definicion}

El límite existe si y solo si existen los dos laterales y coinciden (\cref{eje:2-laterales}). Por ejemplo, para $f(x)=\abs{x}/x$ los laterales en~0 son $1$ y $-1$, y el límite no existe.

\section{Límites notables}

\begin{proposicion}[$\sin x/x$]\label{prop:senx}
Para $0<\abs{x}<\pi/2$ se cumple $\cos x<\dfrac{\sin x}{x}<1$. En consecuencia, $\displaystyle\lim_{x\to0}\frac{\sin x}{x}=1$ y $\displaystyle\lim_{x\to0}\frac{1-\cos x}{x^{2}}=\frac12$.
\end{proposicion}

\begin{proof}
Sea $0<x<\pi/2$ y considérense, en la circunferencia unidad, el triángulo de vértices $O=(0,0)$, $A=(1,0)$ y $P=(\cos x,\sin x)$, el sector circular $OAP$ y el triángulo $OAT$ con $T=(1,\tan x)$. Están contenidos uno en otro, luego sus áreas cumplen $\tfrac12\sin x<\tfrac12x<\tfrac12\tan x$. Dividiendo por $\tfrac12\sin x>0$ y tomando inversos, $\cos x<\sin x/x<1$. Las tres expresiones son pares, luego la desigualdad vale también para $-\pi/2<x<0$. Como $\cos x\to1$ (continuidad del coseno, \cref{tema:3}), el sándwich da $\sin x/x\to1$. Por último,
\[
\frac{1-\cos x}{x^{2}}=\frac{1-\cos^{2}x}{x^{2}(1+\cos x)}=\Bigl(\frac{\sin x}{x}\Bigr)^{2}\frac{1}{1+\cos x}\longrightarrow 1\cdot\frac12.\qedhere
\]
\end{proof}

Para la exponencial se adopta el punto de vista del curso 18.01 del MIT~\cite{mit1801} (sesiones~16 a~19): el número $e$ es la única base $b>1$ para la que la gráfica de $b^{x}$ tiene pendiente~1 en el origen. Se admite, por tanto, que la exponencial es continua y que
\begin{equation}\label{eq:lim-exp}
\lim_{h\to0}\frac{e^{h}-1}{h}=1,
\end{equation}
y que esta base coincide con el límite del \cref{teo:e}. En Spivak~\cite[cap.~18]{spivak} se construyen $\ln$ y $\exp$ a partir de la integral y todo lo anterior se demuestra.

\begin{proposicion}[Límites notables de la exponencial y el logaritmo]\label{prop:notables-exp}
\[
\text{(i)}\ \lim_{x\to0}\frac{\ln(1+x)}{x}=1;\qquad
\text{(ii)}\ \lim_{x\to0}(1+x)^{1/x}=e;\qquad
\text{(iii)}\ \lim_{n\to\infty}\Bigl(1+\frac{t}{n}\Bigr)^{n}=e^{t}\ \ (t\in\R).
\]
\end{proposicion}

\begin{proof}
(i) Sea $u=\ln(1+x)$, que tiende a~0 cuando $x\to0$ (continuidad del logaritmo) y es distinto de~0 si $x\ne0$. Como $x=e^{u}-1$, $\dfrac{\ln(1+x)}{x}=\dfrac{u}{e^{u}-1}\to1$ por~\eqref{eq:lim-exp}. (ii) $(1+x)^{1/x}=\exp\bigl(\ln(1+x)/x\bigr)\to e^{1}$, por (i) y la continuidad de la exponencial. (iii) Para $t=0$ es trivial. Si $t\ne0$, con $x_n=t/n\to0$, $(1+t/n)^{n}=\bigl((1+x_n)^{1/x_n}\bigr)^{t}\to e^{t}$ por (ii) y el \cref{teo:heine}.
\end{proof}

\begin{table}[h]
\centering\small
\caption{Límites notables en $x\to0$. Se demuestran en este tema o en sus ejercicios.}\label{tab:notables}
\begin{tabular}{@{}lll@{}}
\toprule
$\dfrac{\sin x}{x}\to1$ & $\dfrac{1-\cos x}{x^{2}}\to\dfrac12$ & $\dfrac{\tan x}{x}\to1$ \\[2mm]
$\dfrac{e^{x}-1}{x}\to1$ & $\dfrac{\ln(1+x)}{x}\to1$ & $(1+x)^{1/x}\to e$ \\[2mm]
$\dfrac{a^{x}-1}{x}\to\ln a$ & $\dfrac{(1+x)^{p}-1}{x}\to p$ & $\dfrac{\tanh x}{x}\to1$ \\
\bottomrule
\end{tabular}
\end{table}

\section{Asíntotas}

\begin{definicion}[Asíntotas]\label{def:asintotas}
Sea $f$ una función.
\begin{enumerate}
  \item La recta $x=a$ es una \emph{asíntota vertical} si alguno de los límites laterales de $f$ en $a$ es $+\infty$ o $-\infty$.
  \item La recta $y=L$ es una \emph{asíntota horizontal} en $+\infty$ si $\lim_{x\to+\infty}f(x)=L$ (análogamente en $-\infty$).
  \item La recta $y=mx+b$ con $m\ne0$ es una \emph{asíntota oblicua} en $+\infty$ si $\lim_{x\to+\infty}\bigl(f(x)-mx-b\bigr)=0$.
\end{enumerate}
\end{definicion}

\begin{proposicion}[Cálculo de la asíntota oblicua]\label{prop:oblicua}
$y=mx+b$ es asíntota oblicua de $f$ en $+\infty$ si y solo si $m=\lim_{x\to+\infty}f(x)/x$ (finito y no nulo) y $b=\lim_{x\to+\infty}\bigl(f(x)-mx\bigr)$ (finito).
\end{proposicion}

\begin{proof}
Si $f(x)-mx-b\to0$, entonces $f(x)-mx\to b$, y dividiendo por $x$, $f(x)/x-m=(f(x)-mx)/x\to0$. Recíprocamente, si existen ambos límites, $f(x)-mx-b\to b-b=0$.
\end{proof}

\begin{ejemplo}[Asíntotas de una función racional y de la softplus]\label{ej:asintotas}
(a) $f(x)=\dfrac{x^{2}+1}{x-1}=x+1+\dfrac{2}{x-1}$ (división de polinomios). Tiene asíntota vertical $x=1$, con $f\to+\infty$ por la derecha y $f\to-\infty$ por la izquierda, y asíntota oblicua $y=x+1$ en $\pm\infty$, porque $2/(x-1)\to0$.
(b) $\softplus(x)\to0$ cuando $x\to-\infty$ (asíntota horizontal $y=0$) y $\softplus(x)-x=\softplus(-x)\to0$ cuando $x\to+\infty$ (asíntota oblicua $y=x$). Es lo que muestra la \cref{fig:activaciones}: la softplus se pega a la ReLU lejos del origen.
(c) $\sigma$ tiene asíntotas horizontales $y=0$ en $-\infty$ e $y=1$ en $+\infty$.
\end{ejemplo}

\begin{observacion}[Errores frecuentes]\label{obs:2-errores}
(a) En una demostración $\varepsilon$--$\delta$, elegir $\delta$ dependiente de $x$: $\delta$ solo puede depender de $\varepsilon$ (y de $a$).
(b) Escribir $\infty-\infty=0$, $0\cdot\infty=0$ o $1^{\infty}=1$: son \emph{formas indeterminadas}, cuyo valor depende de las funciones concretas; por ejemplo, $(1+1/n)^{n}\to e\ne1$.
(c) Concluir que $\lim_{x\to a}f(x)=f(a)$ sin justificar la continuidad; el límite no depende del valor $f(a)$.
(d) Aplicar el álgebra de límites cuando alguno no existe: $\lim(f+g)$ puede existir aunque no existan $\lim f$ ni $\lim g$.
(e) Creer que una sucesión acotada converge: $(-1)^{n}$ es acotada y divergente; hace falta además la monotonía (\cref{teo:monotona}).
\end{observacion}

\begin{observacion}[Conexión con el curso]
En el módulo~1 se demuestra que la softmax con temperatura tiende al vector indicador del máximo cuando $T\to0^{+}$ y a la distribución uniforme cuando $T\to\infty$; son límites de funciones como los de este tema, y la demostración usa que $e^{-c/T}\to0$ para $c>0$. El \cref{teo:monotona} es la herramienta con la que, en el módulo~7, se prueba que la pérdida decrece y converge a lo largo del descenso de gradiente, y la \cref{prop:notables-exp}(iii) reaparece en el módulo~21 como el factor $(1-1/n)^{n}\to e^{-1}$.
\end{observacion}

\begin{lecturas}
\item Spivak~\cite{spivak}, cap.~5 «Limits» y cap.~22 «Infinite Sequences».
\item Gowers y Chua~\cite{cambA1}, cap.~2 «Convergence of sequences» (§2.1--2.3) y §5.1 «Limits».
\item OpenStax~\cite{osc1}, §2.2 «The Limit of a Function», §2.3 «The Limit Laws», §2.5 «The Precise Definition of a Limit» y §4.6 «Limits at Infinity and Asymptotes».
\item Stewart~\cite{stewart}, §2.4 «The Precise Definition of a Limit» y §2.6 «Limits at Infinity; Horizontal Asymptotes».
\end{lecturas}

\section*{Ejercicios del tema 2}
\addcontentsline{toc}{section}{Ejercicios del tema 2}

\begin{ejercicio}{1}{eje:2-negacion}
Escriba con cuantificadores la negación de «$a_n\to L$» y la de «$\lim_{x\to a}f(x)=L$». Use la primera para explicar qué hay que exhibir para probar que $a_n\not\to L$.
\end{ejercicio}

\begin{ejercicio}{1}{eje:2-N}
Para $a_n=\dfrac{2n}{n+3}$, que tiende a~2, halle el menor $N$ tal que $\abs{a_n-2}<0{,}01$ para todo $n\ge N$. Repita con $\varepsilon=10^{-6}$.
\end{ejercicio}

\begin{ejercicio}{2}{eje:2-sucesiones-calculo}
Calcule los límites: (a) $\dfrac{3n^{2}-n+4}{5n^{2}+2}$; (b) $\sqrt{n^{2}+n}-n$; (c) $\dfrac{2^{n}+3^{n}}{3^{n+1}}$; (d) $\dfrac{\sin n}{n}$; (e) $\Bigl(1-\dfrac{2}{n}\Bigr)^{n}$; (f) $\sqrt[n]{2^{n}+3^{n}}$.
\end{ejercicio}

\begin{ejercicio}{2}{eje:2-funciones-calculo}
Calcule, si existen, los límites siguientes (indique $\pm\infty$ cuando proceda):
\begin{enumerate}[label=(\alph*)]
  \item $\lim_{x\to2}\dfrac{x^{2}-4}{x^{2}-3x+2}$;
  \item $\lim_{x\to1^{+}}\dfrac{x+1}{x^{2}-1}$ y $\lim_{x\to1^{-}}\dfrac{x+1}{x^{2}-1}$;
  \item $\lim_{x\to\infty}\dfrac{2x^{3}-x}{x^{3}+4x^{2}}$;
  \item $\lim_{x\to0}\dfrac{\sqrt{1+x}-1}{x}$;
  \item $\lim_{x\to-2}\dfrac{x^{3}+8}{x+2}$.
\end{enumerate}
\fuente{adaptado de MIT 18.01SC, Jerison y Mattuck~\cite{mit1801ex}, ejercicio 1D-1 (límites de funciones racionales).}
\end{ejercicio}

\begin{ejercicio}{2}{eje:2-notables}
Usando la \cref{tab:notables} (y demostrando de paso las entradas que no se han probado en el texto), calcule: (a) $\lim_{x\to0}\dfrac{\sin 5x}{3x}$; (b) $\lim_{x\to0}\dfrac{\tan x}{x}$; (c) $\lim_{x\to0}\dfrac{e^{2x}-1}{\sin x}$; (d) $\lim_{x\to0}\dfrac{2^{x}-1}{x}$; (e) $\lim_{x\to0}\dfrac{\tanh x}{x}$; (f) $\lim_{x\to\infty}x\sin(1/x)$.
\end{ejercicio}

\begin{ejercicio}{2}{eje:2-asintotas}
Halle todas las asíntotas de: (a) $f(x)=\dfrac{2x^{2}-3x}{x+1}$; (b) $g(x)=\dfrac{x}{\sqrt{x^{2}+1}}$; (c) $h(x)=xe^{1/x}$ (para la oblicua, use la \cref{tab:notables}); (d) $k(x)=\ln(1+e^{2x})$.
\end{ejercicio}

\begin{ejercicio}{3}{eje:2-eN}
Demuestre con la definición $\varepsilon$--$N$: (a) $\dfrac{n^{2}}{n^{2}+1}\to1$; (b) $\dfrac{1}{\sqrt n}\to0$; (c) $\dfrac{(-1)^{n}}{n^{2}}\to0$; (d) $\dfrac{n+5}{2n-1}\to\dfrac12$.
\end{ejercicio}

\begin{ejercicio}{3}{eje:2-ed}
Demuestre con la definición $\varepsilon$--$\delta$: (a) $\lim_{x\to2}(5x+8)=18$; (b) $\lim_{x\to2}(x^{2}+2x)=8$; (c) $\lim_{x\to3}\dfrac{x^{2}-9}{x-3}=6$; (d) $\lim_{x\to0}x^{4}=0$.
\fuente{adaptado de OpenStax~\cite{osc1}, §2.5, ejercicios de la sección.}
\end{ejercicio}

\begin{ejercicio}{3}{eje:2-ed-dificiles}
Demuestre con la definición $\varepsilon$--$\delta$: (a) $\lim_{x\to4}\sqrt{x}=2$ (indicación: $\sqrt x-2=\frac{x-4}{\sqrt x+2}$); (b) $\lim_{x\to1}\dfrac1x=1$; (c) $\lim_{x\to0}\dfrac{1}{x^{2}}=+\infty$, con la \cref{def:limites-variantes}(ii).
\fuente{adaptado de OpenStax~\cite{osc1}, §2.5, ejercicios de la sección (límites infinitos).}
\end{ejercicio}

\begin{ejercicio}{3}{eje:2-no-existe}
Demuestre que no existen: (a) $\lim_{x\to1}\dfrac{\abs{x-1}}{x-1}$, directamente con la definición $\varepsilon$--$\delta$ (tome $\varepsilon=1$); (b) $\lim_{x\to0}f(x)$, donde $f(x)=1$ si $x\in\Q$ y $f(x)=0$ si $x\notin\Q$, con el \cref{teo:heine}.
\fuente{adaptado de OpenStax~\cite{osc1}, §2.5, ejercicios de la sección.}
\end{ejercicio}

\begin{ejercicio}{3}{eje:2-producto}
Sean $\lim_{x\to a}f(x)=L$ y $\lim_{x\to a}g(x)=M$. Demuestre directamente con la definición $\varepsilon$--$\delta$, sin usar sucesiones, que $\lim_{x\to a}f(x)g(x)=LM$. (Indicación: siga el esquema de la demostración del \cref{teo:algebra-sucesiones}(iii) y acote primero $f$ cerca de $a$.)
\fuente{adaptado de OpenStax~\cite{osc1}, §2.5, ejercicio de la sección sobre la ley del producto.}
\end{ejercicio}

\begin{ejercicio}{3}{eje:2-cociente}
(a) Sea $b_n\to M\ne0$ con $b_n\ne0$ para todo $n$. Demuestre que existe $N$ con $\abs{b_n}>\abs{M}/2$ para $n\ge N$, y deduzca que $1/b_n\to1/M$. (b) Sea $(a_n)$ una sucesión de reales no nulos. Para cada una de las condiciones siguientes, demuestre que $(1/a_n)$ converge o dé un contraejemplo: (i) $a_n\to L$ para algún $L\in\R$; (ii) $a_n\to L$ para algún $L\ne0$.
\fuente{adaptado de Cambridge, \emph{Tripos} IA 2019, Paper~1~\cite{tripos2019}, pregunta~3E.}
\end{ejercicio}

\begin{ejercicio}{3}{eje:2-alternada}
(a) Demuestre que $((-1)^{n})$ diverge. (b) Más en general, demuestre que si $a_n\to L$ entonces $a_{n+1}-a_n\to0$, y deduzca que $(r^{n})$ diverge para $r\le-1$. (c) ¿Es cierto el recíproco de (b)? Considere $a_n=\sqrt n$.
\fuente{adaptado de Gowers y Chua~\cite{cambA1}, §2.1 y §2.2.}
\end{ejercicio}

\begin{ejercicio}{3}{eje:2-recurrencia}
Sea $a_1=1$ y $a_{n+1}=\sqrt{2+a_n}$. Demuestre por inducción que $0<a_n<2$ y que $(a_n)$ es creciente. Deduzca que converge y calcule su límite.
\end{ejercicio}

\begin{ejercicio}{3}{eje:2-raizn}
(a) Demuestre que $\sqrt[n]{n}\to1$. (Indicación: escriba $\sqrt[n]{n}=1+h_n$ con $h_n\ge0$ y use $(1+h_n)^{n}\ge\binom n2h_n^{2}$.) (b) Deduzca que $\sqrt[n]{a}\to1$ para todo $a>0$.
\end{ejercicio}

\begin{ejercicio}{3}{eje:2-laterales}
Demuestre que $\lim_{x\to a}f(x)=L$ si y solo si $\lim_{x\to a^{+}}f(x)=L$ y $\lim_{x\to a^{-}}f(x)=L$. Aplíquelo a $f(x)=\ReLU(x)/x$ en $a=0$.
\end{ejercicio}
